\documentclass[10pt,a4paper]{article}
\usepackage[margin=15mm]{geometry}
\usepackage[T1]{fontenc}
\usepackage{tgheros,tgcursor}
\renewcommand{\familydefault}{\sfdefault}
\usepackage{microtype,booktabs,tabularx,array,xcolor}
\usepackage[colorlinks=true,urlcolor=blue!55!black]{hyperref}
\definecolor{ink}{HTML}{16324F}
\definecolor{accent}{HTML}{087F8C}
\definecolor{muted}{HTML}{526273}
\pagestyle{empty}
\setlength{\parindent}{0pt}
\setlength{\parskip}{4pt}
\setlength{\tabcolsep}{5pt}
\renewcommand{\arraystretch}{1.12}
\newcommand{\sect}[1]{\vspace{4pt}{\color{ink}\bfseries #1}\par\vspace{1pt}}
\newcolumntype{Y}{>{\raggedright\arraybackslash}X}
\hypersetup{pdftitle={Four-H100 training improvements},pdfsubject={Measured trainer optimizations and validation}}
\begin{document}
{\LARGE\bfseries\color{ink} Four-H100 training improvements}\hfill{\small 26--27 September 2026}

{\large\bfseries\color{accent} 90.6 iterations/s sustained live training\quad +25.7\%}

Versus 72.1 it/s with the corrected loader and starting GPU kernels. Four H100 SXM 80 GB;
global batch 131,072; FP32; AdamW; no-PSQT threats network, L1=1024 and L2=L3=32.
All production filters remain enabled, including check/capture filtering and random FEN skipping of 10.

\sect{Measured progression}
\begin{tabularx}{\linewidth}{@{}Yrr@{}}
\toprule
\textbf{Mature-network configuration} & \textbf{Live it/s} & \textbf{Gain over baseline}\\
\midrule
Corrected loader, starting GPU kernels & 72.10 & ---\\
FT launch/atomic tuning + bucketed dense L1 + DDP & 85.39 & 18.4\%\\
Add weighted-NNZ L1 forward & 86.40 & 19.8\%\\
Add eight-position FT gradient aggregation & \textbf{90.62} & \textbf{25.7\%}\\
\bottomrule
\end{tabularx}
{\small\color{muted} Medians of three 2,048-step windows after 1,024 warmup steps. Final cached throughput: 93.76 it/s.}

\sect{Retained improvements}
\begin{tabularx}{\linewidth}{@{}p{39mm}Y@{}}
\textbf{Loader allocation reuse} & glibc arena/mmap/trim settings reduce native batch release from 10.5 to 0.024 ms.\\
\textbf{Loader staging threads} & Four torch threads and 64 loader workers per rank; combined loader tuning raises roughly 37 to 72 it/s. Production sampling is preserved.\\
\textbf{BMI2 detection} & Match the complete CPU flag so AVX-512 VBMI2 does not disable the decoder fast path. No isolated speedup attributed.\\
\textbf{Loss metric updates} & Remove per-update GPU-to-CPU NaN checks; retain nonfinite-loss checks at logging boundaries. Included in the corrected baseline.\\
\textbf{FT forward launch} & 256-thread blocks for H100/L1=1024: about 1.69 to 1.50 ms. General shapes retain their existing path.\\
\textbf{FT backward launch} & Four 128-thread column blocks plus zero-gradient atomic guards: about 5.34 to 3.84 ms. These are separate commits.\\
\textbf{FT feature aggregation} & Combine updates across eight positions before global atomics. About 3.18 ms backward + 0.18 ms packing; +4.9\% live throughput. Duplicate indices use a correctness fallback.\\
\textbf{Selected-bucket L1} & Compute the selected bucket directly. Input-gradient kernel: 0.381 to 0.104 ms; weight-gradient kernel: 0.330 to 0.107 ms, plus small reductions.\\
\textbf{Weighted-NNZ L1} & Compact nonzeros and coalesce weight reads. Forward: 0.315 to 0.166 ms including packing/transpose; about +1.3\% whole-step throughput at 25\% activation density.\\
\textbf{DDP bucket size} & 50 to 400 MiB reduces gradient collectives from three to one. Pre-aggregation four-GPU scaling was 3.18$\times$ (79.6\% efficiency); scaling is not proven optimal.\\
\end{tabularx}

\sect{Hardware-counter finding}
A full 5090 VM enabled Nsight Compute counters. With cache flushing disabled, FT forward reached
\textbf{94\% L2 / 26\% DRAM} throughput; direct backward reached \textbf{73\% L2 / 12\% DRAM}.
This points to cache/atomic pressure on that GPU. Aggregation reduced L2 reduction sectors by 17.7\%
in the default-cache run. These are \textbf{5090 diagnostics, not H100 bandwidth measurements}.

\sect{Validation and limits}
\textbf{104 tests passed, 2 skipped on H100}; full production-size validation and checkpoint resume
through step 2,560 completed with finite loss. The mature net came from the supplied public
\texttt{nn-252f33942263.nnue}; comparisons used identical fresh optimizer state.
Sparse L1 can regress dense early nets by about 1\%. Quantization fusion, separate bias reductions,
and forced NCCL protocols did not yield useful whole-training gains. No strength claim is made.
After rebase, each branch passes 47 local CPU tests (59 skipped). GPU timings predate the
history rewrite; measured runtime sources are preserved.

\sect{Branches and evidence}
\begin{tabularx}{\linewidth}{@{}p{62mm}Y@{}}
\texttt{optimize/4xh100-no-loader} & GPU/DDP work, baseline loader launch settings.\\
\texttt{optimize/4xh100-with-loader} & Adds separate BMI2, allocation-reuse, and staging-thread commits.\\
\end{tabularx}
Both are based on master \texttt{aa2fff2}; upstream HyperLogLog changes remain in both.
Added Markdown reports remain untracked and unstaged.

{\small\raggedright Evidence: \path{docs/benchmarks/h100-followup-2026-09-26/}.\\Workload:
\href{https://github.com/vondele/nettest/blob/9bc8d4f4f89b9126079176dfdcd901d7bf294f7d/threats.yaml}{nettest threats.yaml}.
Source: \texttt{h100-improvements.tex}.}
\end{document}
